\section{The result and conventions}\label{s:result}

We study the spin-\(2\) square-lattice AKLT model of Affleck, Kennedy,
Lieb, and Tasaki~\cite{AKLT88,KLT88}. We prove uniqueness of its ground
state on every torus with integer periods \(L,N\ge3\) and a uniform
positive spectral gap for all integers \(L,N\ge L_0\), with no parity
restriction, for one fixed threshold \(L_0\). The computational inputs
are local operator inequalities, an open-trail growth bound, and a
type-dependent positive-vector bound obtained from stationary Bellman
inequalities and finite prefix sums. All computations are checked in
exact arithmetic. The polynomial realization of the spin space,
including its Hilbert space normalization, is fixed in
Definition~\ref{d:poly}.

\begin{definition}\label{d:model}
For integers \(L,N\ge3\), let \(\Lambda=\Lambda_{L,N}=\mathbb Z_L\times\mathbb Z_N\)
with its simple, unoriented nearest-neighbor edge set \(E_{L,N}\).
The Hilbert space is \(\mathcal H_\Lambda=\bigotimes_{x\in\Lambda}V_4\).
Each site carries the spin-\(2\) space \(V_4\), and \(P_{\{x,y\}}\)
is the orthogonal projection onto total spin \(4\) in \(V_4\otimes V_4\).
Fix one orientation of each edge and set
\[
 q_{xy}=z_{x,1}z_{y,2}-z_{x,2}z_{y,1},\qquad
 B_{L,N}=\prod_{e\in E_{L,N}}q_e,\qquad
 H_{L,N}=\sum_{e\in E_{L,N}}P_e.
\]
Every edge is counted once. The spectral gap of \(H_{L,N}\) means its
smallest positive eigenvalue.
\end{definition}

Here an edge joins two distinct sites whose difference is one of
\((1,0),(-1,0),(0,1),(0,-1)\), with coordinates reduced modulo the
periods. The edge set is a set of unordered pairs, so that an edge
is included only once even if it has more than one description.
The notation \(q_e\) uses the chosen orientation; \(P_e\) does not.
The polynomial \(B_{L,N}\) is the valence-bond state.

\begin{theorem}\label{t:main}
For every \(L,N\ge3\), \(B_{L,N}\ne0\) and
\(\ker H_{L,N}=\mathbb C B_{L,N}\).
There exist \(c>0\) and an integer \(L_0\ge3\) such that, for all
integers \(L,N\ge L_0\), every nonzero eigenvalue of \(H_{L,N}\)
is at least \(c\).
\end{theorem}

\begin{proof}
Proposition~\ref{p:toruskernel} proves the nonvanishing and kernel
assertions for every \(L,N\ge3\). For the uniform gap,
Theorem~\ref{t:inputs} supplies (H1), (H2) with
\[
 C_V=10\ge1,\qquad
 0<q=\frac{1119948287353}{1120437614225}<1,
\]
and (H3\(_h\)) for the prescribed type-dependent kernel and fixed
orderings. The common thresholds are \(c_{\rm s}=961/1000\) and
\(c_{\rm t}=187/200\).
These are exactly the hypotheses of Theorem~\ref{t:rec}, which
gives finite \(A\ge2\) and \(0<\theta<1\) such that
\(\mathsf P(A,\theta)\) holds.
Theorem~\ref{t:cutMOS} then gives
\(\mathsf W(64,p_0,M_\star)\) for the plane frame family, with
one \(p_0\in(0,1]\) and one integer \(M_\star\ge1\) valid for
every proper plane three-coloring and every translate.
This is the hypothesis of the periodic assembly theorem,
Theorem~\ref{t:mainW}, proved in Section~4, for the plane frame
family of Definition~\ref{d:framemodel}. Its direct quantum
conclusion supplies \(c>0\) and an integer \(L_0\ge3\) with the
asserted spectral bound for all integers \(L,N\ge L_0\).
\end{proof}

\subsection{Structure of the proof}

The uniqueness assertion has a separate algebraic proof. The kernel
of each edge projection consists of multiples of its edge determinant.
Unique factorization combines these divisibilities, and degree four
at every torus site leaves a scalar quotient. This argument applies
also when either period is three.

For the gap, finite fourth moments in Section~2 give tight frames of
coherent states. In Section~3 we assign three explicit rotated
icosahedral frames according to a proper torus coloring. Their
separation makes the squared modulus of the frame image of
\(B_{L,N}\) strictly positive. The resulting antiferromagnetic
probability law on the torus has a single-site heat bath whose gap
bounds the quantum gap from below. A conditional physical square
lifts to the plane, where a coordinatewise sign change converts its
law to the positive-interaction plane specification. This sign
change is defined on the lift.

Section~4 uses a uniform estimate for window marginals to choose one
square scale. A stopped coupling bounds the Hamming cost of changing
one exterior letter. Counting all periodic translates of that square
gives a gap for block updates, and comparison of conditional block
and site Dirichlet forms gives the torus site gap.
Theorem~\ref{t:mainW} combines this with the quantum comparison of
Theorem~\ref{t:QD} to obtain the uniform quantum gap.
Section~5 proves the required window estimate from the pair-response
bound \(\mathsf P(A,\theta)\), uniformly over plane colorings and
placements. Sections~6--8 prove that response bound using a walk
expansion, a contracting positive-vector recursion, and the exact
computational inputs. Thus the gap implication is
\[
 \begin{gathered}
 \text{exact computational inputs}
 \ \Longrightarrow\ \mathsf P
 \ \Longrightarrow\ \mathsf W\\
 \ \Longrightarrow\ \text{periodic block gap}
 \ \Longrightarrow\ \text{site gap}
 \ \Longrightarrow\ \text{quantum gap}.
 \end{gathered}
\]
The finite probability and comparison tools are proved in Appendix~A;
the planar topology and the computational certificate arguments are
proved in Appendices~B--D. All scale choices precede the choice of the
torus periods, coloring, block placement and exterior configuration.

\emph{Formal verification.}
Theorem~\ref{t:main} has been formalized in Lean~4 and checked by the
Lean kernel. The formalization has two parts. The first derives
Theorem~\ref{t:main} from (H1) and (H3\(_h\)) and proves (H2); it is
self-contained and is checked in minutes. The second checks the
finite inputs (H1) and (H3\(_h\)) in the Lean kernel, which takes
about \(152\) and \(23\) CPU-hours. Both parts are available with
this paper. % CHECK: add repository reference

\subsection{Conventions}

All Hilbert spaces are complex, with inner products conjugate linear
in the first argument. An inequality between self-adjoint operators
means the corresponding inequality of quadratic forms.
We abbreviate \(B=B_{L,N}\) and \(H=H_{L,N}\) when the periods are
fixed. Products over empty sets are one and sums over empty sets
are zero.

On the plane,
\[
 d(x,y)=|x_1-y_1|+|x_2-y_2|,\qquad
 |x-y|_\infty=\max\{|x_1-y_1|,|x_2-y_2|\}.
\]
We write \(x\sim y\) for nearest neighbors and
\(d(x,S)=\min_{y\in S}d(x,y)\), with the value \(+\infty\) if
\(S=\varnothing\). Distances in plane estimates are these ambient
distances, even if sites have been removed. A torus distance, when
needed, is stated explicitly.
Lattice translations lie in \(\mathbb Z^2\), and a radius used as
an endpoint of a lattice interval is a nonnegative integer.

For probabilities \(\nu,\nu'\) on a finite set,
\[
 d_{\rm TV}(\nu,\nu')
 =\frac12\sum_s|\nu(s)-\nu'(s)|
 =\max_A|\nu(A)-\nu'(A)|.
\]
For two configurations on the same product of site alphabets, their
Hamming distance is the number of sites at which their letters
differ. Conditional expectations and probabilities are finite sums
on positive conditioning events. Variances use squared complex
moduli. A claimed Poincare gap \(g>0\) means
\(\mathcal D(f)\ge g\operatorname{Var}(f)\) for every \(f\).
This convention includes a one-point configuration space, for which
both sides vanish; the infimum over positive-variance functions is
then \(+\infty\). Every heat-bath update rate is one per physical
site unless a different set of update labels is explicitly specified.
